small
Die folgenden Themen werden in den T0-Dokumenten ausführlich behandelt. Die Dokumentnummern verweisen auf das vollständige Verzeichnis im nächsten Anhang.

\section{Kosmologie ohne Expansion}
Statisches, ewig existierendes Universum ohne Urknall (\href{https://github.com/jpascher/T0-Time-Mass-Duality/blob/main/2/pdf/025_T0_Kosmologie_De.pdf}{Dok.~025}, 026, 156). Rotverschiebung als geometrischer Effekt im $\xsi$-Feld: $z \approx \xsi \cdot \ln(d/\lP)$. Hubble-Konstante $H_0 \approx 66{,}2$~km/s/Mpc aus T0 abgeleitet, was die Hubble-Spannung löst (\href{https://github.com/jpascher/T0-Time-Mass-Duality/blob/main/2/pdf/026_T0_Geometrische_Kosmologie_De.pdf}{Dok.~026}). CMB und Casimir-Effekt als zwei Manifestationen desselben $\xsi$-Feldes (\href{https://github.com/jpascher/T0-Time-Mass-Duality/blob/main/2/pdf/025_T0_Kosmologie_De.pdf}{Dok.~025}, 091). CMB-Dipol-Anomalie erklärt durch wellenlängenabhängige Rotverschiebung (\href{https://github.com/jpascher/T0-Time-Mass-Duality/blob/main/2/pdf/039_Zwei-Dipole-CMB_De.pdf}{Dok.~039}). Dunkle Energie als geometrische Illusion, Dunkle Materie als Torsions-Haltefaktor (\href{https://github.com/jpascher/T0-Time-Mass-Duality/blob/main/2/pdf/149_FFGFT-torsion_De.pdf}{Dok.~149}).

\section{Teilchenmassen und Massenformeln}
Vollständige Herleitung aller Leptonmassen aus $\xsi$ und $f = 7500$ (\href{https://github.com/jpascher/T0-Time-Mass-Duality/blob/main/2/pdf/006_T0_Teilchenmassen_De.pdf}{Dok.~006}, 046, 149). Quarkmassen geometrisch abgeleitet (\href{https://github.com/jpascher/T0-Time-Mass-Duality/blob/main/2/pdf/149_FFGFT-torsion_De.pdf}{Dok.~149}). Koide-Formel als mathematische Konsequenz von $\xsi$ bewiesen (\href{https://github.com/jpascher/T0-Time-Mass-Duality/blob/main/2/pdf/116_T0_koide-formel-3_De.pdf}{Dok.~116}). Neutrino-Masse $m_\nu = \xsi_0^2/2 \times m_e = 4{,}54$~meV (\href{https://github.com/jpascher/T0-Time-Mass-Duality/blob/main/2/pdf/007_T0_Neutrinos_De.pdf}{Dok.~007}, 047). Higgs-VEV $v \approx 246{,}7$~GeV mit $0{,}2\%$ Genauigkeit und $W$/$Z$-Boson-Massen aus Torsionsgeometrie (\href{https://github.com/jpascher/T0-Time-Mass-Duality/blob/main/2/pdf/149_FFGFT-torsion_De.pdf}{Dok.~149}).

\section{Anomale magnetische Momente}
Elektron-, Myon- und Tau-Anomalie rein geometrisch hergeleitet (\href{https://github.com/jpascher/T0-Time-Mass-Duality/blob/main/2/pdf/018_T0_Anomale-g2-10_De.pdf}{Dok.~018}, 149). Hausdorff-Dimensionen $p \approx 5/3$ (Myon) und $4/3$ (Tau) für fraktale Verzweigungen. Tau-Anomalie-Vorhersage $a_\tau \approx 1{,}282 \times 10^{-3}$, testbar bei Belle~II (\href{https://github.com/jpascher/T0-Time-Mass-Duality/blob/main/2/pdf/149_FFGFT-torsion_De.pdf}{Dok.~149}).

\section{Lagrangian-Formulierungen}
Vereinfachter T0-Lagrangian $\mathcal{L}_0 = \varepsilon(\partial\Delta m)^2$ (\href{https://github.com/jpascher/T0-Time-Mass-Duality/blob/main/2/pdf/019_T0_lagrndian_De.pdf}{Dok.~019}, 095). Erweiterter Lagrangian mit Selbstwechselwirkung (\href{https://github.com/jpascher/T0-Time-Mass-Duality/blob/main/2/pdf/019_T0_lagrndian_De.pdf}{Dok.~019}). Vergleich mit Standard-QFT (\href{https://github.com/jpascher/T0-Time-Mass-Duality/blob/main/2/pdf/049_LagrandianVergleich_De.pdf}{Dok.~049}). Notwendigkeit zweier Lagrangians (\href{https://github.com/jpascher/T0-Time-Mass-Duality/blob/main/2/pdf/095_Notwendigkeit_zwei_lagrange_De.pdf}{Dok.~095}). Dirac-Gleichung in T0 (\href{https://github.com/jpascher/T0-Time-Mass-Duality/blob/main/2/pdf/050_diracVereinfacht_De.pdf}{Dok.~050}, 051, 053).

\section{Elimination der Masse als fundamentale Größe}
$E = mc^2$ ist mathematisch identisch mit $E = m$; $c$ ist nur ein Verhältnis (\href{https://github.com/jpascher/T0-Time-Mass-Duality/blob/main/2/pdf/077_E-mc2_De.pdf}{Dok.~077}). Systematische Elimination von Masse aus allen Grundgleichungen (\href{https://github.com/jpascher/T0-Time-Mass-Duality/blob/main/2/pdf/052_EliminationOfMass_De.pdf}{Dok.~052}, 053, 054). Verhältnisbasiertes vs.\ parameterbasiertes Zahlensystem (\href{https://github.com/jpascher/T0-Time-Mass-Duality/blob/main/2/pdf/057_RelokativesZahlensystem_De.pdf}{Dok.~057}).

\section{Bell-Ungleichungen und Nichtlokalität}
T0-modifizierte CHSH-Ungleichung mit $\xsi$-Korrektur (\href{https://github.com/jpascher/T0-Time-Mass-Duality/blob/main/2/pdf/023_Bell_De.pdf}{Dok.~023}, 023a). Lokalität und Realismus wiederhergestellt durch Zeitfeld als lokale versteckte Variable. Superdeterminismus durch gemeinsame Zeitfeld-Geschichte (\href{https://github.com/jpascher/T0-Time-Mass-Duality/blob/main/2/pdf/020_T0_QM-QFT-RT_De.pdf}{Dok.~020}). Scheinbare Instantanität als lokale Constraint-Erfüllung $T \cdot E = 1$ (\href{https://github.com/jpascher/T0-Time-Mass-Duality/blob/main/2/pdf/131_scheinbar_instantan_De.pdf}{Dok.~131}). Dynamische Photonenmasse $m_\gamma = \omega$ (\href{https://github.com/jpascher/T0-Time-Mass-Duality/blob/main/2/pdf/055_DynMassePhotonenNichtlokal_De.pdf}{Dok.~055}).

\section{Bewusstsein, Cortex, DNA}
No-Go-Theorem für Agentität in unitärer QM; T0 ermöglicht Agentität durch fraktale Inkohärenz (\href{https://github.com/jpascher/T0-Time-Mass-Duality/blob/main/2/pdf/100_Consciousness_De.pdf}{Dok.~100}). Cortex-Windungen als biologisches Analogon der 4D-Torsionsstruktur --- 9 Parallelen (\href{https://github.com/jpascher/T0-Time-Mass-Duality/blob/main/2/pdf/154_Cortex_De.pdf}{Dok.~154}). DNA-Doppelhelix als biologische Lösung desselben geometrischen Problems --- 10 Parallelen (\href{https://github.com/jpascher/T0-Time-Mass-Duality/blob/main/2/pdf/155_DNA_De.pdf}{Dok.~155}). Geometrischer Kompatibilismus (\href{https://github.com/jpascher/T0-Time-Mass-Duality/blob/main/2/pdf/100_Consciousness_De.pdf}{Dok.~100}).

\section{Vergleiche mit anderen Theorien}
T0 vs.\ Standardmodell (\href{https://github.com/jpascher/T0-Time-Mass-Duality/blob/main/2/pdf/068_T0vsESM_ConceptualAnalysis_De.pdf}{Dok.~068}), Synergetics (\href{https://github.com/jpascher/T0-Time-Mass-Duality/blob/main/2/pdf/033_T0-Theory-vs-Synergetics_De.pdf}{Dok.~033}), Penrose (\href{https://github.com/jpascher/T0-Time-Mass-Duality/blob/main/2/pdf/030_T0_penrose_De.pdf}{Dok.~030}), Matsas et al.\ (\href{https://github.com/jpascher/T0-Time-Mass-Duality/blob/main/2/pdf/105_Matsas_T0_Vergleich_De.pdf}{Dok.~105}), Donoghue (\href{https://github.com/jpascher/T0-Time-Mass-Duality/blob/main/2/pdf/140_T0_CMB_Donoghue_Analyse_De.pdf}{Dok.~140}), Plasma-Kosmologie/Unnikrishnan (\href{https://github.com/jpascher/T0-Time-Mass-Duality/blob/main/2/pdf/036_T0_peratt_De.pdf}{Dok.~036}).

\section{Experimentelle Tests und Technologie}
Photonenchip-Optimierung durch T0-Prinzipien (\href{https://github.com/jpascher/T0-Time-Mass-Duality/blob/main/2/pdf/083_T0_photonenchip-china_De.pdf}{Dok.~083}, 084, 085). RSA-Kryptographie und Quantencomputer-Benchmarks (\href{https://github.com/jpascher/T0-Time-Mass-Duality/blob/main/2/pdf/075_RSA_De.pdf}{Dok.~075}, 076). Atominterferometrie, LIGO/Virgo-Korrekturen, 73-Qubit-Bell-Tests (\href{https://github.com/jpascher/T0-Time-Mass-Duality/blob/main/2/pdf/020_T0_QM-QFT-RT_De.pdf}{Dok.~020}, 147). Turing-Muster und biologische Selbstorganisation (\href{https://github.com/jpascher/T0-Time-Mass-Duality/blob/main/2/pdf/146_turing_De.pdf}{Dok.~146}).

\section{Mathematische Struktur und Ontologie}
Zirkularität der Konstanten (\href{https://github.com/jpascher/T0-Time-Mass-Duality/blob/main/2/pdf/101_zirkularitaet-Konstanten_De.pdf}{Dok.~101}). Ontologische Hierarchie (\href{https://github.com/jpascher/T0-Time-Mass-Duality/blob/main/2/pdf/152_ontologische-ord_De.pdf}{Dok.~152}). Energie-Reduktion auf 5 ontologische Ebenen (\href{https://github.com/jpascher/T0-Time-Mass-Duality/blob/main/2/pdf/153_energie-reduktion-on_De.pdf}{Dok.~153}). Netzwerkdarstellung des T0-Modells als multidimensionale Informationsstruktur (\href{https://github.com/jpascher/T0-Time-Mass-Duality/blob/main/2/pdf/024_T0_netze_De.pdf}{Dok.~024}). Warum Zahlenverhältnisse nicht gekürzt werden dürfen (\href{https://github.com/jpascher/T0-Time-Mass-Duality/blob/main/2/pdf/070_Mathematische_struktur_De.pdf}{Dok.~070}). Kompatibilität 4D-Torsion/fraktal (\href{https://github.com/jpascher/T0-Time-Mass-Duality/blob/main/2/pdf/150_kompatiblitaet_De.pdf}{Dok.~150}). Musikalische Spirale: 137 als Resonanzpunkt, $(4/3)^{137} \approx 2^{57}$ (\href{https://github.com/jpascher/T0-Time-Mass-Duality/blob/main/2/pdf/060_musical-spiral-137-_De.pdf}{Dok.~060}). Was \emph{ist} das Universum? Ein universelles Energiefeld $E(x,t)$ mit $\Box E = 0$ und $\xsi$ (\href{https://github.com/jpascher/T0-Time-Mass-Duality/blob/main/2/pdf/156_was-ist-universum_De.pdf}{Dok.~156}).
